% s09-compat.tex —— §7 Local/global compatibility，原文 p.34（§7 标题）–p.36
\section{局部/整体相容}
\label{sec:compat}

我们这里要谈的所有结果都推广到 Hodge 型 Shimura 簇（[2]）。
为简化叙述，我们宁愿限制在 PEL 情形\yizhuw{PEL 型 Shimura 数据指带极化（Polarization）、自同态（Endomorphism）与 level 结构的数据；此处按惯例保留缩写。}。

设 $(H,X)$ 是 PEL 型的 Shimura 数据。
设 $H_{\QQ_p}=G$，即前面那个拟分裂 $p$-adic 群，其中 $E=\QQ_p$。
我们还设 $H_{\QQ_p}$ 非分歧。
固定一个 $p$ 之外的充分小的 level $K^p$、一个嵌入 $\overline{\QQ}\hookrightarrow\overline{\QQ}_p$，并记 $E$ 为反射域相应的 $p$-adic 完备化（是的，这里有记号冲突……但另一个 $E$ 现在就是 $\QQ_p$）。
我们在 $G(\QQ_p)$ 中固定一个超特殊紧子群。
记 $\underline{G}$ 为 $G$ 在 $\ZZ_p$ 上相应的约化模型。
对紧开子群 $K_p\subset\underline{G}(\ZZ_p)$，记
\[
\mathrm{Sh}_{K_p}
\]
为 $E$ 上的刚性解析空间，它是水平为 $K^pK_p$ 的 Shimura 簇的解析化中万有阿贝尔概形的良好归约轨迹。
若 $K_p$ 是我们的超特殊子群，\srcpageinl{35}则这就是 Kottwitz 整模型之 $p$-adic 完备化的一般纤维。
记 $\mathcal{S}$ 为该整模型，$\widehat{\mathcal{S}}_\eta=\mathrm{Sh}_{\underline{G}(\ZZ_p)}$。

回忆 Scholze 的以下定理（这是重写过的版本，Scholze 的定理更强，但这已是我们唯一需要的东西）。

\begin{theorem}[{[28]}]
pro-étale 层 $\varprojlim_{K_p}\mathrm{Sh}^\diamond_{K_p}$ 可由一个 $E$-完满空间 $\mathrm{Sh}_\infty$ 表示。
此外，这个完满空间表示大解析位点上如下预层所对应的层：该预层把 $(R,R^+)$ 映为四元组 $(A,\lambda,\iota,\bar\eta^p)\in\mathcal{S}(R^+)$（[19]），连同 $A[p^\infty]\otimes_{R^+}R$ 的一个与其 $G$-结构相容的平凡化。
\end{theorem}

我们现在使用 §6 的结果。
由于前一定理的最后一个断言，我们只需对 $S(R^+)$ 的元素（其中 $(R,R^+)$ 是 $E$ 上的 affinoid 完满空间）定义随后出现的所有态射。
任何这样的元素都给出 $R^+$ 上的一个 $p$-可除群。

这个 $p$-可除群的协变 $F$-等晶体（连同其 $G$-结构），连同其 Hodge 滤过，定义一个态射
\[
g:\mathrm{Sh}_\infty\longrightarrow\mathcal{BT}\xrightarrow{\ u\ }\mathrm{Hecke}^\mu,
\]
其中 $\mu$ 由整体 Shimura 数据经从 $\QQ$ 到 $\QQ_p$ 的转移\yizhuw{原文此处作“transfert”，是法语拼法，英文应为 transfer；指由 $\QQ$ 上的 Shimura 数据转移到 $\QQ_p$ 上所得的 $\mu$。}得到。
为此我们使用推论 6.10（PEL 型 level 结构的纳入是直接的）。
这里我们改动了预叠 $\mathcal{BT}$ 的定义 6.8，以纳入 $G$-结构。
于是我们有交换图
\[
\begin{tikzcd}[row sep=3.6em, column sep=2.6em]
 & & \mathcal{F}\ell_\mu^\diamond \arrow[d, "i"'] \arrow[drr] & & \\
\mathrm{Sh}_\infty^\flat \arrow[urr, "\pi_{HT}^\diamond"] \arrow[rr, "g"] \arrow[dr] & & \mathrm{Hecke}^\mu \arrow[dl, "\overleftarrow{h}"'] \arrow[dr, "\overrightarrow{h}"] & \square & \Spa(E)^\diamond \arrow[dl, "{(x_1,can)}"'] \\
 & \Bun_G & & \Bun_G\times\Spa(\QQ_p)^\diamond &
\end{tikzcd}
\]
其中：
\begin{itemize}
\item $\mathcal{F}\ell_\mu=\Gr^{\le\mu}=\Gr^\mu$ 是 $E$ 上联系于 $\mu$ 的旗簇；
\item $\pi_{HT}$ 是 Hodge--Tate 周期态射（这个图实际上给出了 $\pi_{HT}$ 的一个定义）；
\item 态射 $(x_1,can)$ 由 $x_1:\Spa(\FF_p)\to\Bun_G$（由平凡向量丛诱导）与联系于 $E|\QQ_p$ 的 $can$ 给出；
\item 右方块是笛卡尔的；
\item $g$ 经 $\pi_{HT}$ 的分解，由万有阿贝尔概形的 Tate 模在 $p$ 处的无穷 level 结构给出：后者给出相应的 $G$-丛 $V_p(\widetilde{H})$ 的一个平凡化（记号同 6.10），从而给出经由右方块所给的纤维积的分解。
\end{itemize}

局部/整体相容条件于是表述如下。

\begin{quote}
\noindent\textbf{局部/整体相容}：\itshape 设 $\Pi$ 是 $H$ 的自守表示，且 $\Pi_p$ 离散。
则存在重数 $m\in\mathbb{N}$，使得
\[
R\pi_{HT*}\bigl(\overline{\QQ}_\ell[\dim\mathrm{Sh}]\bigr)[\Pi^p]
=m\cdot i^*\overleftarrow{h}^*\mathscr{F}_{\varphi_{\Pi_p}},
\]
其中左端是 Caraiani--Scholze 反常层（[2]。\yizhuw{原文此处仅作“([2].”，缺右括号，疑为“([2])”之误；译文照录原文的不完整括号。}
\end{quote}

\begin{remark}
复合
\[
\mathcal{F}\ell^\diamond_\mu\xrightarrow{\ i\ }\mathrm{Hecke}^\mu
\xrightarrow{\ \overleftarrow{h}\ }\Bun_G
\]
通过拉回 $\Bun_G$ 的 HN-分层\yizhuw{HN 即 Harder--Narasimhan：按向量丛的斜率序列所作的分层，见 §2.4。}，定义了 $\mathcal{F}\ell_\mu$ 的一个分层。
这就是所谓 $\mathcal{F}\ell_\mu$ 的 Newton 分层，见 [2]。
\end{remark}

\srcpage{36}
\begin{remark}[尖点性条件的动机]
这里给出一个来自 [2] 的、对尖点性条件的动机。
在 [2] 中，作者们证明了 $R\pi_{HT*}\bigl(\overline{\QQ}_\ell[\dim\mathrm{Sh}]\bigr)[\Pi^p]$ 沿 Newton 分层局部常，其纤维是相应 Igusa 簇的上同调\yizhuw{Igusa 簇是 Newton 分层各层上的覆盖，专门用来在整模型上“看见”该层的部分上同调，见 [12]、[20]。}（[12]，[20]）。
这是在完满设定下对 Mantovan 在 [20] 中工作的重新解释，其中邻近循环 $R\Psi_{\bar\eta}$ 被 $R\pi_{HT*}$ 取代。
利用 Shin 关于 Igusa 簇上同调的交代和的公式（[31]，推广了 [12]），可以看到这个 Igusa 簇的上同调的交代和，作为 $J_b(E)$ 的表示，是用一个从 $G(\QQ_p)$ 的表示到 $M_b(\QQ_p)$ 的表示的 Jacquet 函子计算的\yizhuw{Jacquet 函子沿抛物子群取其幂零根基下的 Jacquet 模（余不变量），从而把 $G$ 的表示送到 Levi 子群的表示。}，其中 $M_b$ 是作为斜率态射的中心化子的 Levi 子群。
因此若 $b$ 不是基本的，这个交代和为零。
与局部/整体相容条件相结合，这构成尖点性条件的一个强动机。
\end{remark}

\begin{remark}[等特征情形]
当然，在等特征情形，用 Shtuka 的模空间应当有同类的局部/整体相容性；此时的好处是 $\mu$ 不必是 minuscule 的，这扩大了局部/整体相容可能适用的范围。
\end{remark}

\begin{remark}[$\ell=p$]
$p$-adic Banach 层 $R\pi_{HT*}\widehat{\mathcal{O}}_{\mathrm{Sh}_\infty}$ 及其挠版本 $R\pi_{HT*}\mathcal{O}^+_{\mathrm{Sh}_\infty}/p$ 是有意义的。
从这个观点看，可以问我们的局部猜想在所谓 p-adic Langlands 纲领的框架中是否有一个 $p$-adic 对应物。
若确乎如此，则对非基本的 $b$，$\widetilde{J}_b$ 在 $x_b^*\mathscr{F}_\varphi$ 上的作用将不再经 $\pi_0(\widetilde{J}_b)=J_b(E)$ 的作用分解（见注 4.6），这应当导致有趣的现象。
例如对 $\mathrm{GL}_2/\QQ_p$ 与模曲线、对 ordinary 的 $b$，有 $\widetilde{J}_b=\mathbb{B}^{\varphi=p}$（见例 2.24），人们可以希望把这与 $p$-adic 分布的空间联系起来。
\end{remark}
